
\begin{abstract}

We recover the finite matching extensions and regular halved hypertopes associated with every centrally symmetric spherical regular polytope. One orbital-presentation theorem and one parity rewriting give the full defining presentations. The exceptional calculation is an exact orbit census in five finite Coxeter actions. Four published censuses are correct. For the 120-cell there are 35 unordered diagonal classes, rather than the asserted 15; the published presentation omits 20 orbit relations and defines an infinite group for every parameter. We give the missing relations explicitly. The intended finite 120-cell construction, its order and its halved geometry remain valid with the corrected presentation. The original construction and families are attributed to their sources; the correction is distinct from a new open-conjecture solution.

\end{abstract}

\section{Source contract and common notation}\label{source-contract-and-common-notation}

The source is C.A. Piedade, \emph{Infinite Families of Hypertopes from Centrally Symmetric Polytopes}, Electronic Journal of Combinatorics \textbf{30}(2) (2023), P2.20, DOI 10.37236/10392, final 21-page journal edition. The targets are the general orbital presentation (Corollary 6), all families in §§3--4, including their full presentations, orders, geometric type and stated toroidal residues, and the iteration observation in §5. Classical background and the retained McMullen--Schulte extension and toroidal-coordinate results remain explicit inputs. We do not claim a new classification of arbitrary hypertopes or a historical priority determination.

Write \(L=[p_1,\ldots,p_{n-1}]=\langle t_0,\ldots,t_{n-1}\rangle\) for the spherical group of \(P\). In the extension put
\[
 a=\rho_0,\quad b=\rho_1=t_0,\quad c_j=\rho_j=t_{j-1}\ (2\le j\le n),
 \quad d=aba.
\]
Let \(\mathcal F(p_1,\ldots,p_{n-1})\) denote the fork Coxeter presentation on \(d,b,c_2,\ldots,c_n\): both \(dc_2\) and \(bc_2\) have label \(p_1\), the remaining chain has labels \(p_2,\ldots,p_{n-1}\), and all undisplayed pairs commute. When \(n=2\), this is the string \(d,c_2,b\) with labels \(p_1,p_1\).

The companion paper \emph{Orbit presentations and parity kernels of matching extensions} is denoted U.M. Its Theorems 1--3 give complete presentations and actual regular geometries. We use \(\theta\) for substitution \(b\mapsto d\), fixing every \(c_j\). For any old-group word \(w\), an additional relation in the extension and its halved form are respectively
\[
 (a w^{-1}a w)^e=1,
 \qquad (\theta(w)^{-1}w)^e=1.                 \tag{1}
\]
For a Coxeter word \(\beta=b c_2\cdots c_n\), put \(\widetilde\beta=d c_2\cdots c_n\). If \(\alpha=\beta^h\) is the antipodal involution, its relation is
\[
 (a\beta^h a\beta^h)^s=1,
 \qquad (\widetilde\beta^h\beta^h)^s=1.       \tag{2}
\]
All statements below hold for every integer \(s\ge2\).

\section{The finite family theorem}\label{the-finite-family-theorem}

\textbf{Theorem 1.} For each input in Table 1, the matching extension \(E_s(P)\) is a regular lattice polytope of type \(\{4,p_1,\ldots,p_{n-1}\}\), with full group \(D_s^q\rtimes L\). Its halving is a regular hypertope with diagram \(\mathcal F(p_1,\ldots,p_{n-1})\), full type-preserving group \(\ker(D_s^q\to C_2)\rtimes L\), and order in the last column. The defining presentations are precisely those of §§3--6, including the additional 20 relations in the last case.

\begin{longtable}{@{}lrrr@{}}
\caption{The finite families.}\tabularnewline
\toprule\noalign{}
Input \(P\) & \(q\) & \(|L|\) & Halved order \\
\midrule\noalign{}
\endfirsthead
\toprule\noalign{}
Input \(P\) & \(q\) & \(|L|\) & Halved order \\
\midrule\noalign{}
\endhead
\bottomrule\noalign{}
\endlastfoot
\(2p\)-gon, \(p\ge2\) & \(p\) & \(4p\) & \(2p(2s)^p\) \\
\(n\)-orthoplex, \(n\ge3\) & \(n\) & \(2^n n!\) & \((4s)^{n-1}(2s)n!\) \\
\(n\)-cube, \(n\ge3\) & \(2^{n-1}\) & \(2^n n!\) & \((2s)^{2^{n-1}}2^{n-1}n!\) \\
Icosahedron \(\{3,5\}\) & \(6\) & \(120\) & \(60(2s)^6\) \\
Dodecahedron \(\{5,3\}\) & \(10\) & \(120\) & \(60(2s)^{10}\) \\
24-cell \(\{3,4,3\}\) & \(12\) & \(1152\) & \(576(2s)^{12}\) \\
600-cell \(\{3,3,5\}\) & \(60\) & \(14400\) & \(7200(2s)^{60}\) \\
120-cell \(\{5,3,3\}\) & \(300\) & \(14400\) & \(7200(2s)^{300}\) \\
\end{longtable}

\textbf{Proof.} The classical classification of finite irreducible string Coxeter groups supplies exactly these centrally symmetric convex polytopes: the remaining simplex groups are centerless in the relevant ranks. All input face posets are lattices. The numbers of vertices and the group orders are the usual spherical Coxeter data; for the exceptional rows they are also verified in the exact actions of §5. Apply U.M (2), U.M Theorem 3 and U.M (5). Completeness of the specialized presentations is proved below. Since the order is strictly increasing in \(s\), each row is an infinite family of pairwise nonisomorphic finite geometries. $\square$

The finite Coxeter classification and the extension construction are retained classical results. The source's order and existence claims in the 120-cell row refer to this well-defined finite construction. They do not make its incomplete displayed presentation finite.

\section{Polygons and orthoplexes}\label{polygons-and-orthoplexes}

The following two applications use the common presentation and parity theorems. Their arguments reuse the previously completed polygon and orthoplex proofs. Their conventional toroidal lattice names use McMullen--Schulte Corollaries 8C7 and 8C6 and §6D, in exactly the form quoted in source Corollaries 9 and 12.

For a \(2p\)-gon, retain its extended presentation on \(a,b,c=c_2\) of type \([4,2p]\), with
\[
 [ab(cb)^{j-1}]^4=1\ (2\le j\le p-1),\qquad
 [ab(cb)^{p-1}]^{2s}=1.                       \tag{3}
\]
This is also the hole presentation \(\{4,2p\mid4^{p-2},2s\}\) in the source's Definition 8. The superscript \(p-2\) here denotes a row of \(p-2\) labels, not exponentiation of a relation. Empty index ranges at \(p=2\) impose no relation.

Put \(w_j=b(cb)^{j-1}\). U.M Theorem 2 gives
\[
 R_j=\theta(w_j)w_j=(dc)^{j-1}db(cb)^{j-1}.
\]
Thus the full halved presentation consists of \([2p,2p]\) on \(d,c,b\), together with \(R_j^2=1\) for \(2\le j\le p-1\) and \(R_p^s=1\). This proves source Proposition 11, including sufficiency, rather than merely checking its words in a finite group. The geometry has a string diagram and is therefore a regular map of type \(\{2p,2p\}\). At \(p=2\), the retained original toroid is \(\{4,4\}_{(2s,0)}\). Realize its universal square grid on $\mathbb Z^2$, with translation lattice
\[
\Lambda=\langle(2s,0),(0,2s)\rangle.
\]
The black vertices are $B=\{(x,y):x+y\equiv0\pmod2\}$. In the halved map, whose string order is $d,c,b$, an old square is an edge joining its two black corners, and a white vertex becomes the face bounded by its four black neighbors. Thus the black grid has edge basis $u=(1,1)$ and $v=(1,-1)$. The integral coordinate isomorphism
\[
 B\longrightarrow\mathbb Z^2,\qquad
 (x,y)\longmapsto\left(\frac{x+y}{2},\frac{x-y}{2}\right)
\]
carries $u,v$ to the standard basis and carries the two generators of $\Lambda$ to $(s,s)$ and $(s,-s)$. These generate the square-lattice period group customarily denoted by $(s,s)$; changing the sign of its second generator gives the equivalent quarter-turn basis $(s,s),(-s,s)$. Because $\Lambda\subset B$, this coordinate map descends bijectively on vertices, and the described edges and faces show that it preserves the complete square-grid incidence. The halving is therefore exactly \(\{4,4\}_{(s,s)}\), including its translation lattice, not merely a map of the same order.

For the \(n\)-orthoplex, retain the classical cubical toroid
\[
 E_s(P)=\{4,3^{n-2},4\}_{(2s,0^{n-1})}.
\]
Its complete extra relation over the affine string is \((aw)^{2s}=1\), where
\[
 t=c_2\cdots c_n c_{n-1}\cdots c_2,\qquad w=btb.
\]
Since \(b,d\) commute,
\[
 \theta(w)w=dtdbtb=(dtb)^2.
\]
The full halved presentation is consequently \(\mathcal F(3^{n-2},4)\) together with
\[
 (d c_2\cdots c_n c_{n-1}\cdots c_2 b)^{2s}=1. \tag{4}
\]
The original translation lattice lies in the parity kernel: its stated generators are conjugates of a word containing \(2s\) copies of \(a\). Thus the halving is the affine fork reflection geometry modulo the same translation lattice, with the source's notation \(\{3/3,3^{n-3},4\}_{(2s,0^{n-1})}\). Its actual maximal residues are spherical by U.M §4, since every proper orthoplex face is a simplex. Its diagram is affine; it is locally spherical of Euclidean type. This proves Proposition 13 in its all-rank scope, with no finite-test extrapolation.

\section{Cubes and their toroidal residues}\label{cubes-and-their-toroidal-residues}

Let \(P\) be the \(n\)-cube, \(n\ge3\). Use vertices \((\pm1)^n\); \(t_0\) negates the first coordinate and \(t_j\) interchanges coordinates \(j,j+1\). Starting at \((1,\ldots,1)\), the Coxeter element \(\beta=t_0\cdots t_{n-1}\) produces a vertex with exactly \(i\) negative coordinates for \(1\le i\le n\). The stabilizer of the base vertex is the full coordinate-permutation group. It is transitive on vertices with any given number of negatives. Hamming distance is preserved by the full signed-permutation group, so two such classes cannot merge even when a pair is reversed. There are exactly \(n\) unordered diagonal classes, represented by \(\beta^i\), and \(\beta^n\) is the antipodal involution. This proves source Lemma 14 without retaining its separate assertion about the number of classes.

The first class is already supplied by the initial Coxeter label four. U.M Theorems 1--2 therefore give the complete additional relations
\[
 (a\beta^{-i}a\beta^i)^2=1\quad(2\le i<n),
 \qquad (a\beta^n a\beta^n)^s=1,              \tag{5}
\]
for \(E_s(P)\), and
\[
 (\widetilde\beta^{-i}\beta^i)^2=1\quad(2\le i<n),
 \qquad (\widetilde\beta^n\beta^n)^s=1        \tag{6}
\]
for the halving. Their Coxeter parts are \([4,4,3^{n-2}]\) and \(\mathcal F(4,3^{n-2})\), respectively.

Every square face of the cube contains no antipodal pair. U.M (6) identifies its extended lower face with \(2^{\{4\}}\), with all four independent commuting vertex flips and its dihedral square action. This is exactly the polygon application with \(p=s=2\), hence the actual marked map \(\{4,4\}_{(4,0)}\). Its halving is exactly \(\{4,4\}_{(2,2)}\). This proves the residue assertions in Corollary 15 and Proposition 16; the argument identifies the marked face, rather than drawing a conclusion from a relation giving only a quotient.

For \(n=3\), all other maximal residues are spherical: deleting a new vertex gives the cube, and deleting the branch node gives a product of rank-one spherical geometries. Thus the halving is locally toroidal in the source's rank-four definition. Equations (5)--(6) include the source's additional parameter relation \((\widetilde\beta^3\beta^3)^s=1\); the toroidal face relation follows from the identified lower face. For larger \(n\) we assert the specified toroidal rank-three residues, not a definition of locally toroidal hypertopes in every rank. The incidental source comparison with the universal locally toroidal catalogue on p.14 is not needed for, and is outside, this principal family contract.

\section{The exceptional orbit census and four unchanged presentations}\label{the-exceptional-orbit-census-and-four-unchanged-presentations}

For each exceptional \(L=[p_1,\ldots,p_{n-1}]\), the attached GAP program starts from its literal finite Coxeter presentation. It enumerates cosets of \(L_0=\langle t_1,\ldots,t_{n-1}\rangle\), obtains the distinguished permutations, and checks their exact group order, point stabilizer and central antipodal involution \(\beta^h\). It then computes all orbits on unordered pairs of vertices, not only pairs reached by powers of \(\beta\). Every pair is included and the orbits partition the full pair set. A separate Burnside calculation over every group element uses the fixed-pair formula \(\binom{f(g)}2+c_2(g)\), where \(f(g)\) counts fixed vertices and \(c_2(g)\) counts two-cycles. It confirms each orbit count, including the 120-cell numerator \(504000=35\cdot14400\). The following results are checked against the declared row data.

\begin{longtable}{@{}lrrrp{57mm}@{}}
\caption{Exceptional representative exponents.}\tabularnewline
\toprule\noalign{}
Input & Vertices & \(|L_0|\) & \(h\) & Complete representative exponents \\
\midrule\noalign{}
\endfirsthead
\toprule\noalign{}
Input & Vertices & \(|L_0|\) & \(h\) & Complete representative exponents \\
\midrule\noalign{}
\endhead
\bottomrule\noalign{}
\endlastfoot
Icosahedron & 12 & 10 & 5 & \(1,3,5\) \\
Dodecahedron & 20 & 6 & 5 & \(1,2,3,4,5\) \\
24-cell & 24 & 48 & 6 & \(1,3,4,6\) \\
600-cell & 120 & 120 & 15 & \(1,4,6,7,9,10,12,15\) \\
\end{longtable}

The corresponding orbit-size lists are respectively
\[
 (30,30,6),\quad (30,60,60,30,10),\quad
 (96,72,96,12),\quad
 (720,1200,720,1800,720,1200,720,60).
\]
Their sums are \(\binom{|V|}{2}\). Every displayed representative belongs to a distinct orbit; the last orbit is exactly the antipodal matching. For the first row, powers \(1,2\) coincide and powers \(3,4\) coincide; for the 24-cell, powers \(1,2\) coincide and powers \(4,5\) coincide. These observations also check the corresponding source explanations. The use of a finite Coxeter presentation and its classical order is explicit: the computed permutation image has that full order, so it is the actual spherical action, not a smaller unexplained quotient.

\textbf{Theorem 2.} For each of these four inputs let \(J\) be its list of representative exponents with \(1\) and \(h\) removed. The extension has the complete presentation with its string Coxeter part and relations
\[
 (a\beta^{-i}a\beta^i)^2=1\ (i\in J),\qquad
 (a\beta^h a\beta^h)^s=1.                    \tag{7}
\]
The halving has its fork Coxeter part and the complete additional relations
\[
 (\widetilde\beta^{-i}\beta^i)^2=1\ (i\in J),\qquad
 (\widetilde\beta^h\beta^h)^s=1.             \tag{8}
\]

\textbf{Proof.} The finite census proves exactly the orbit-representative hypothesis of U.M Theorem 1. Apply that theorem and then the complete parity rewriting. The orders, full groups, geometry and marking were proved in Theorem 1. $\square$

This recovers Corollaries 17,19,21,23 and Propositions 18,20,22,24. The icosahedron and 600-cell have only simplex proper faces, so U.M §4 proves actual spherical maximal residues for every \(s\). The two fork diagrams are the standard compact hyperbolic diagrams; we retain their classical Coxeter recognition, as in the earlier generalized-cube application. At \(s=2\), the groups are the old even-coordinate Danzer groups \(2^{11}\rtimes[3,5]\) and \(2^{119}\rtimes[3,3,5]\), respectively. These are the source's previously known examples, not additional new constructions.

For the 24-cell, an octahedral facet has six vertices and no antipodal pair. Its actual lower face is \(2^{\{3,4\}}\). The orthoplex application at \(n=3,s=2\) identifies it with \(\{4,3,4\}_{(4,0,0)}\), and identifies its halving with \(\{3/3,4\}_{(4,0,0)}\). Thus the claimed quotients in Corollary 21 and Proposition 22 are in fact these marked toroids themselves. All other maximal residues of the halving are spherical, as follows from its lower and upper face sections: the vertex residues are the 24-cell, the remaining proper lower faces are simplices. We state this exact residue property; the source explicitly leaves the general terminology in ranks above four undeveloped.

\section{The corrected 120-cell presentation}\label{the-corrected-120-cell-presentation}

Here \(L=[5,3,3]\), \(\beta=t_0t_1t_2t_3\), and \(\alpha=\beta^{15}\). The vertex stabilizer is \([3,3]\), of order24; there are 600 vertices. The point-stabilizer action has 45 orbits including the fixed base vertex. Its nontrivial ordered-pair orbitals merge, on forgetting orientation, into \textbf{35} unordered-pair orbits. Powers \(\beta^i\), \(1\le i\le15\), represent 15 distinct such orbits. They do not exhaust them.

The twenty missing orbits have representatives \(w=\beta^a t_j\beta^b\) in Table 3. The action and word convention are the ones fixed in §1; indices \(j\) start at zero. The orbit-size column counts unordered pairs, not vertices in a stabilizer orbit.

\begin{longtable}[]{@{}rrrr@{}}
\caption{The twenty omitted 120-cell pair orbits.}\tabularnewline
\toprule\noalign{}
\(a\) & \(j\) & \(b\) & Orbit size \\
\midrule\noalign{}
\endfirsthead
\toprule\noalign{}
\(a\) & \(j\) & \(b\) & Orbit size \\
\midrule\noalign{}
\endhead
\bottomrule\noalign{}
\endlastfoot
\(-3\) & 3 & \(-1\) & 3600 \\
\(-4\) & 0 & \(-2\) & 1200 \\
\(-4\) & 3 & \(-1\) & 7200 \\
\(-5\) & 0 & \(-2\) & 2400 \\
\(-5\) & 3 & \(-1\) & 7200 \\
\(-4\) & 3 & \(-2\) & 3600 \\
\(-6\) & 3 & \(-1\) & 7200 \\
\(-9\) & 1 & 4 & 1200 \\
\(-5\) & 3 & \(-2\) & 7200 \\
\(-7\) & 3 & \(-1\) & 14400 \\
\(-5\) & 3 & \(-3\) & 1800 \\
\(-8\) & 3 & \(-1\) & 7200 \\
\(-7\) & 3 & \(-2\) & 7200 \\
\(-6\) & 3 & \(-4\) & 1200 \\
\(-9\) & 3 & \(-1\) & 3600 \\
\(-8\) & 3 & \(-2\) & 7200 \\
\(-7\) & 3 & \(-4\) & 2400 \\
\(-10\) & 3 & \(-1\) & 7200 \\
\(-7\) & 3 & \(-5\) & 1200 \\
\(-10\) & 3 & \(-2\) & 3600 \\
\end{longtable}

The complete pair-size multiset is
\[
 300^1,\ 1200^6,\ 1800^1,\ 2400^2,\ 3600^6,\ 7200^{18},\ 14400^1,
                                                               \tag{9}
\]
where the exponents denote multiplicities. These 35 orbits contain all 179700 unordered pairs. The literal attached census verifies each short word, distinctness from all fifteen old classes, mutual distinctness, and full coverage. In particular the correction is an explicit presentation independent of any GAP coset numbering.

\textbf{Theorem 3 (corrected last family).} The intended finite extension has complete Coxeter part \([4,5,3,3]\), the old relations
\[
 (a\beta^{-i}a\beta^i)^2=1\ (2\le i\le14),\qquad
 (a\beta^{15}a\beta^{15})^s=1,
\]
and the additional twenty relations
\[
 (a w^{-1}a w)^2=1\qquad(w\text{ in Table 3}). \tag{10}
\]
Its halving has complete Coxeter part \(\mathcal F(5,3,3)\), the old halved relations for \(2\le i\le14\) and the antipodal relation, and the additional twenty relations
\[
 (\theta(w)^{-1}w)^2=1\qquad(w\text{ in Table 3}). \tag{11}
\]
For \(w=\beta^a t_j\beta^b\), the word \(\theta(w)\) is explicitly \(\widetilde\beta^a\theta(t_j)\widetilde\beta^b\), where \(\theta(t_0)=d\) and \(\theta(t_j)=c_{j+1}\) for \(j\ge1\). The full group orders remain \(14400(2s)^{300}\) and \(7200(2s)^{300}\).

\textbf{Proof.} Table 3 together with the fifteen old powers is a complete set of unordered-pair representatives. Apply U.M Theorems 1--3. All nonantipodal labels are now two, so its Coxeter kernel is exactly \(D_s^{300}\). This supplies sufficiency of both corrected presentations, all group assertions, and the actual regular halved geometry. $\square$

\textbf{Theorem 4 (literal source presentation).} The presentation printed in Corollary 25, and the fork presentation printed in Proposition 26, define infinite groups for every integer \(s\ge2\).

\textbf{Proof.} The old fifteen representatives are distinct, with the last one antipodal. By U.M Theorem 1, the literal Corollary 25 presentation is \(W(m)\rtimes[5,3,3]\), where the 20 missing orbit classes have label \(\infty\), the other nonantipodal classes label two, and the antipodal class label \(s\). For instance \(w=\beta^{-3}t_3\beta^{-1}\), the first row of Table 3, is omitted. U.M Theorem 4 shows that \(u_vu_{v^w}\) has infinite order. It already lies in the parity kernel. The complete Schreier calculation identifies that kernel with the literal Proposition 26 presentation, proving the second assertion as well. $\square$

This error is not resolved by verifying the printed words in the finite group: that finite group is a quotient of the infinite presented group. The statement that the finite construction satisfies those words is true; their claimed sufficiency as defining relations is false. No claim about the source's general construction is refuted.

\section{Iteration, literal corrections and scope}\label{iteration-literal-corrections-and-scope}

The product of central dihedral half-turns in U.M §5 gives a fixed-point-free central involution of every extension in Table 1 when \(s\) is even. The retained lattice property persists. Iterating the construction therefore gives types \(\{4,4,p_1,\ldots,p_{n-1}\}\), and the endpoint-halving theorem applies again. For the source's suggested toroidal inputs \(\{4,4\}_{(2s,0)}\) and \(\{4,3^{n-2},4\}_{(2s,0^{n-1})}\), with even \(s\ge4\), these are the polygon/orthoplex extensions already established above; the same central symmetry proof and lattice input make them eligible. We do not infer lattice preservation for arbitrary graph covers.

The complete final primary was read. Physical pages18--20 were also visually checked for the following literal discrepancies. The Table 2 24-cell row prints \((\widetilde\beta^6\beta^{15})^s\), whereas Proposition 22 and the antipodal calculation require \((\widetilde\beta^6\beta^6)^s\). The Table 2 120-cell row omits the equality separating its two types of relations; even reading it with the equality supplied leaves the substantive missing-orbit error proved above. Corollary 25's continuation labels its group by the dual input \(\{3,3,5\}\) instead of \(\{5,3,3\}\). Equation (3) on p.8 gives one excessive generator index in its case range; the actual extension has \(n+1\) generators. These typographical issues are separate from the complete-presentation correction.

The source also discusses earlier catalogues and the novelty of its families. We do not recover those bibliographical comparisons as mathematical classification statements. A bounded search on 11 October 2026 for the exact title with ``erratum'' and ``correction'', and for Piedade/120-cell/35, did not locate an accessible primary correction. The author's public publication list was checked. This is not an exhaustive priority search, and we make no first-discovery claim.

\section{Proof and certificate boundaries}\label{proof-and-certificate-boundaries}

The companion U.M supplies the uniform arguments. The finite census is exhaustive for the five exceptional finite actions and binds the explicit missing words. It supplies no finite extrapolation in \(s\), \(n\) or \(p\): those parameters are handled by the written normal-form, coordinate and geometric proofs. The auxiliary Lean declarations verify the conjugation identities used in rewriting, the correct inverse-conjugate relation, and the exact powers and infinite order of the omitted-pair matrix block. The passage from the source presentation to a Coxeter semidirect product, the real reflection representation, the retained extension theorem, and the full hypertope proof remain written mathematics.

The separate source ledger records every numbered source statement and the main unnumbered claims. Corollaries 9 and 12 retain the classical coordinate identification of the toroids; the application proves the halved presentations and their unchanged translation lattices. The bipartite case of source Corollary 3 needed by every family is proved in U.M Lemma 3a; its broader general statement remains attributed background. All source family results are therefore recovered or explicitly corrected relative to these named retained inputs. This is a candidate complete principal contract, subject to independent review, not a claim that every sentence of the source has been rederived. No compression ratio or open-question credit is assigned here.

\begin{thebibliography}{9}
\bibitem{P} C. A. Piedade, Infinite Families of Hypertopes from Centrally Symmetric Polytopes, \emph{Electronic Journal of Combinatorics} 30(2) (2023), P2.20, doi:10.37236/10392.
\bibitem{MS} P. McMullen and E. Schulte, \emph{Abstract Regular Polytopes}, Encyclopedia of Mathematics and its Applications 92, Cambridge University Press, 2002. The precise retained results are quoted in \cite{P}; independent access to the book is not claimed.
\bibitem{MW} A. Montero and A. I. Weiss, Locally spherical hypertopes from generalised cubes, \emph{The Art of Discrete and Applied Mathematics} (2020). The halving and simplex-face ideas are attributed to this work; the necessary proofs are written here.
\end{thebibliography}
